\subsubsection{Logical Layers}

The application's logical architecture is visually divided into the following main layers:	

\textbf{Presentation layer}

The Presentation layer corresponds to the frontend component and is responsible for handling the user interface, user experience, and API communication. Each navigation state or view in the application corresponds to a distinct page, and each reusable component can be used across multiple pages.

The user interface is built using React with TypeScript. This strategic choice allows developers to leverage React's component-based architecture, enabling the creation of reusable UI blocks (components) that compose larger screens (pages). Although this approach requires a more complex initial implementation effort, it speeds up development and improves code maintainability and extensibility over time. TypeScript also enforces contract definitions, ensuring consistency across the frontend data flow. 

\textbf{API layer}

The API layer exposes the application's use cases to external parties via RESTful HTTP endpoints. Different controllers handle distinct functional areas (set of use cases), with the main controllers dedicated to Expenses and Reports. Controllers receive HTTP requests, validate them, translate them into commands or queries for the Application layer, and return an HTTP response. This layer, along with the subsequent ones, constitutes the backend component and uses the ASP.NET Core platform.

\vspace{1cm}

\textbf{Application Layer}

The Application layer contains the application logic and is responsible for handling all use cases. This layer orchestrates entities defined in the Domain layer to execute business rules and interacts with repository abstractions (interfaces) to retrieve and persist data. The concrete implementations of these repositories are provided by the Infrastructure layer.


Additionally, this layer defines Data Transfer Objects (DTOs) and interfaces. DTOs standardize data exchange between components, while interfaces are abstractions that specify how external dependencies are accessed without exposing implementation details.

\textbf{Domain layer}

The Domain layer is the core of the application and represents the business model. It encapsulates entities (e.g., \texttt{Expense}, \texttt{User}), value objects, and the most critical and stable business rules. This is a pure layer, with no dependencies on other layers, and it remains independent of how data is stored or presented to the user. This independence ensures a stable and long-lasting business model.

\textbf{Infrastructure layer}

While the Application layer depends on abstractions to perform data access and communicate with external services, the Infrastructure layer provides the corresponding concrete implementations. It is responsible for handling communication with external systems and services. \texttt{ExpenseRepository}, \texttt{BlobStorageService}, and \texttt{DocumentIntelligenceService} are some of the implementations contained in this layer.

This layer is expected to evolve as external technologies and infrastructure requirements change. For example, switching from Blob Storage to another storage provider only requires changes within this layer, without affecting the upper layers, which depend on abstractions rather than concrete implementations.

Another key advantage of using the .NET platform is availability of Entity Framework Core (EF Core) and the built-in Dependency Injection (DI) container. EF Core enables higher-level interaction with the database, allowing data manipulation, relationship management, and migrations through clear, easy-to-understand C\texttt{\#} code. DI is explained in more detail in the next section.
